\documentclass[12pt]{article}
\usepackage[margin=0.65in]{geometry}

\newcommand{\duedate}{02/24/2026}
\newcommand{\assignment}{Representation Learning for Continual Learning}

\newcommand{\name}{Aksel Kretsinger-Walters, adk2164}
\newcommand{\email}{adk2164@columbia.edu}

\makeatletter
\def\input@path{{../}{../../}{../../../}}
\makeatother
\input{pset_template_CLMM.tex} %% DO NOT CHANGE THIS LINE

\begin{document}
\psetheader %% DO NOT CHANGE THIS LINE

% \section*{Representational Continuity for Unsupervised Continual Learning}

\paragraph{Motivation.}
Most continual learning (CL) methods assume supervised labels and focus on preserving task-specific decision
boundaries. However, real-world data streams are often unlabeled and non-stationary. This paper studies
\emph{unsupervised continual learning (UCL)}, where the objective is to learn feature representations
sequentially without catastrophic forgetting.

\paragraph{Core Hypothesis.}
Unsupervised representations are inherently more robust to forgetting because they learn general-purpose
features rather than task-specific discriminative boundaries.

\paragraph{Methodology.}
The authors adapt modern self-supervised learning methods—\textbf{SimSiam} and \textbf{Barlow Twins}—to the
continual learning setting. Representations are trained sequentially across tasks and evaluated using a
frozen encoder with a KNN classifier.

They extend several supervised CL strategies to UCL:
\begin{itemize}
		\item Regularization (Synaptic Intelligence)
		\item Architectural expansion (Progressive Neural Networks)
		\item Rehearsal (Dark Experience Replay)
\end{itemize}

They additionally propose \textbf{Lifelong Unsupervised Mixup (LUMP)}, which interpolates current-task
examples with replay-buffer samples:
\[
		\tilde{x}_{i,\tau} = \lambda x_{i,\tau} + (1-\lambda)x_{j,M}
\]
This encourages linear behavior across tasks and mitigates representational drift.

\paragraph{Key Findings.}
\begin{itemize}
		\item UCL significantly reduces catastrophic forgetting compared to supervised CL.
		\item Simple finetuning in UCL often outperforms sophisticated supervised CL methods.
		\item UCL generalizes better to out-of-distribution (OOD) tasks.
		\item UCL is more robust in few-shot settings.
		\item Loss landscapes of UCL are flatter and smoother than supervised CL.
\end{itemize}

\paragraph{Representational Analysis.}
Using CKA similarity and $\ell_2$ parameter distance:
\begin{itemize}
		\item Two independent UCL models are more similar to each other than two supervised models.
		\item Lower layers remain similar across methods; higher layers diverge.
		\item UCL representations produce coherent, perceptually meaningful feature maps even for earlier tasks.
\end{itemize}

\paragraph{Conclusion.}
Unsupervised continual learning yields more stable, transferable representations than supervised CL. Learning
invariant, general-purpose features implicitly regularizes the learning dynamics and reduces forgetting.

\end{document}
